\documentstyle[%


righttag%


,11pt%


,amssymb%


,russian%               remove in Eng.


,izvru%                 Set izven% in Eng.


]{amsart}





%507 return individually \ch to \cosh, \sh to \sh, \th (!) to \th


\newcommand{\la}{\langle}


\newcommand{\ra}{\rangle}


\newcommand{\tg}{\operatorname{tg}}


\newcommand{\ch}{\operatorname{ch}}


\newcommand{\sh}{\operatorname{sh}}


\newcommand{\th}{\operatorname{th}}


\newcommand{\const}{\operatorname{const}}


\newcommand{\PR}{\operatorname{Pr}}


\newcommand{\tts}[1]{}





%\hoffset=-15mm%                  For Eng.


\setcounter{page}{39}


\god{2001}


\nomer{3~(466)} %                        Remove in Eng.


%\nomer{Vol.\,45, No.\,3}%               Set in Eng.


%\str{\pageref{begin}--\pageref{end}}%  Set in Eng.


\udk{514.753}





\author{\it Л.A.\,Масальцев}


% NB:   ^^ remove in Eng.


\title{Поверхности с плоскими линиями кривизны\\


в пространстве Лобачевского}





\date{}


%\pagestyle{myheadings}%         Set in Eng.


\pagestyle{plain} %              Remove in Eng.


\begin{document}


%\thispagestyle{plain}%          Set in Eng.


\maketitle


\label{begin}





В евклидовом пространстве $R^3$ известно несколько классов поверхностей


с плоскими линиями кривизны: поверхности вращения, резные поверхности,


поверхности Иоахимсталя и другие. В данной статье предлагается


конструкция поверхностей в пространстве Лобачевского $H^3$, аналогичных


по своим свойствам поверхностям Иоахимсталя, т.\,е. дано описание класса


поверхностей, одно семейство линий кривизны которых лежит во


вполне-геодезических плоскостях, содержащих общую геодезическую


простpанства $H^3$. Затем рассмотрены поверхности постоянной средней


кривизны без омбилических точек с одним семейством плоских линий


кривизны. Если минимальная поверхность в $R^3$ имеет одно семейство


плоских линий кривизны, то и другое семейство также является плоским. В


евклидовом пространстве такими поверхностями являются катеноид,


минимальная поверхность Эннепера и однопараметрическое семейство


поверхностей Бонне ([1], c.\,544). Г.\,Венте ([2], теорема 5.5, c.\,54)


доказал, что в пространстве Лобачевского в классе минимальных поверхностей


подобным свойством ``планарности'' одного семейства плоских линий


кривизны обладают только поверхности вращения. В данной статье


показано, что в классе поверхностей постоянной средней кривизны в $H^3$


без омбилических точек только поверхности вращения (гиперболические


поверхности Делоне) имеют одно семейство плоских линий кривизны.


\medskip





{\bf 1.} Рассмотрим модель геометрии Лобачевского в интерпретации Пуанкаре в


верхнем полупространстве $R^3_+$ с метрикой


$ds^2=\frac{dx^2+dy^2+dz^2}{z^2}$. Определим поверхность Иоахимсталя в


$H^3$ как поверхность, у которой каждая линия кривизны из одного


семейства лежит в некоторой вполне-геодезической плоскости, причем все


данные плоскости проходят через фиксированную прямую в $H^3$. Возьмем в


качестве такой оси прямую $l(0,0,z)$. Как известно


([1], c.\,512; [3], c\,.378), в


евклидовом пространстве поверхности Иоахимсталя состоят из множества


ортогональных траекторий семейства сфер переменного радиуса с центрами


на фиксированной оси $l$. По аналогии с евклидовым случаем решим


сначала задачу о нахождении двупараметрического множества ортогональных


траекторий к семейству сфер с центрами на фиксированной оси, а затем


выделим из него класс поверхностей Иоахимсталя. Как известно


([4], c.\,37),


сфера в пространстве Лобачевского с центром в точке $(0,0,s)$


гиперболического радиуса $R(s)$ имеет уравнение


$$


x^2+y^2+(z-s \ch R(s))^2=s^2 \sh^2R(s).


$$





\begin{lem}


Двупараметрическое семейство ортогональных траекторий к


множеству сфер переменного радиуса $R(s)$ с центрами $(0,0,s)$ на оси


$l$ имеет вид


\begin{gather*}


x(s)=\frac{1}{\ch \tau (s)\sqrt {1+c_1^2}},\quad


y(s)=\frac{c_1}{\ch \tau (s)\sqrt {1+c_1^2}},\\


z(s)=-\th \tau (s),


\ \; \text{где}\ \; \tau (s)= \int{\dfrac{(s \ch R(s))'}{s \sh R(s)}ds}


+c_2\quad (c_1,\ c_2 \text{ ---}\const).


\end{gather*}


\end{lem}





{\bf Доказательство.} Обозначим через $(x(s),y(s),z(s))$ координаты


единичного вектора, направленного из точки $(0,0,s \ch R(s))$ в точку


пересечения ортогональной траектории со сферой радиуса $R(s)$.


Уравнение ортогональной траектории можно записать в виде


$$


\ov r(s)=(s \sh R(s)x(s),s \sh R(s)y(s),s \sh R(s)z(s)+s \ch R(s)),


$$


где $x^2(s)+y^2(s)+z^2(s)=1$.


Поскольку данная модель $H^3$ конформно-евклидова, то отыскание


ортогональной траектории приводит к уравнению $\frac{d\ov r}{ds}


=\lambda (s)(x,y,z)$, где $\lambda (s)$ --- некоторый ненулевой множитель.


Отсюда получается система обыкновенных дифференциальных уравнений для


неизвестных функций $x(s)$, $y(s)$, $z(s)$


$$


\frac{(s \sh R(s)x)'}{x}=\frac{(s \sh R(s)y)'}{y}=


\frac{(s \sh R(s)z+s \ch R(s))'}{z}.


$$


Из первого уравнения следует, что $y=c_1x$, где $c_1$ --- некоторая


постоянная. Геометрически это означает, что искомая траектория лежит во


вполне-геодезической плоскости, содержащей ось~$l$. Поскольку


$x^2+y^2+z^2=1$, то отсюда определяем


\begin{equation}


x^2= \frac{1-z^2}{1+c_1^2}.


\end{equation}





Рассмотрим уравнение


\begin{equation}


\frac{(s \sh Rx)'}{x}= \frac{(s \sh Rz+s \ch R)'}{z}.


\end{equation}


После упрощения получим


$$


\frac{x's \sh R}{x}= \frac{s \sh Rz'+(s \ch R)'}{z}.


$$





Продифференцируем уравнение (1) и разделим результат на $2x^2$. Получим


$$


\frac{x'}{x}=-\frac{zz'}{1-z^2}.


$$


Подставим в уравнение (2) и придем к уравнению на функцию $z(s)$


$$


\frac{z'}{z^2-1}=\frac{(s \ch R)'}{s \sh R}.


$$


Интегрируя его (с учетом того, что $|z|<1$), получим


$$


z(s)=- \th \tau (s),


\ \; \text{где}\ \; \tau (s)=\int {\frac{(s \ch R(s))'}{s \sh R(s)}ds}


+c_2.


$$


Затем находим


$$


x(s)=\frac{1}{\ch \tau \sqrt{1+c_1^2}}, \quad


y(s)=\frac{c_1}{\ch \tau \sqrt{1+c_1^2}}.\quad\square


$$





Найденное множество ортогональных траекторий зависит от двух произвольных


постоянных $c_1$, $c_2$. Если связать их между собой, положив $c_1=\tg t$,


$c_2=T(t)$, где $T(t)$ --- произвольная функция от $t$, то получим следующие


уравнения поверхности Иоахимсталя в пространстве Лобачевского $H^3$:


\begin{equation}


\begin{split}


X(s,t)&=\frac{s \sh R(s) \cos t}{\ch \tau (s,t)},\\


Y(s,t)&=\frac{s \sh R(s) \sin t}{\ch \tau (s,t)},\\


Z(s,t)&=s \ch R(s)-s \sh R(s) \th \tau (s,t),\\


\text{где} \; \ \tau(s,t)&=\int


{\frac{(s \ch R(s))'}{s \sh R(s)}ds} +T(t),


\end{split}


\end{equation}


причем $R(s)>0$, $T(t)$ --- произвольные функции переменных $s$ и $t$.





Действительно, вычисляя коэффициенты фундаментальных форм данной поверхности,


убедимся в том, что координатные линии $s$ и $t$ являются линиями кривизны.





\begin{lem}


$\mathrm a)$ Коэффициенты первой квадратичной формы поверхности $(3)$ в


$H^3$ имеют вид


\begin{align*}


g_{11}&=\frac{X_s^2+Y_s^2+Z_s^2}{Z^2}=\bigg(\frac{(s\sh R)'_s-


(s\ch R)'_s\th \tau}{s\ch R-s\sh R\th \tau} \bigg)^2,\\


g_{12}&=\frac{X_sX_t+Y_sY_t+Z_sZ_t}{Z^2}=0,\\


g_{22}&=\frac{X_t^2+Y_t^2+Z_t^2}{Z^2}=(1+T'{}_t^2)\bigg (\frac{s\sh R}


{\ch \tau (s\ch R-s\sh R\th \tau}\bigg)^2;


\end{align*}





$\mathrm b)$ коэффициент $b_{12}$ второй квадратичной формы равен нулю.


\end{lem}





\begin{pf}


Для частных производных вектор-функции $\ov r(s,t)=(X,Y,Z)$ получаем


\begin{alignat*}{2}


X_s&=\frac{\cos t}{\ch \tau}((s\sh R)_s'-\th \tau (s \ch R)_s'),


&\quad X_t&=-\frac{s\sh R}{\ch \tau}(\sin t+\th \tau T_t'),\\


Y_s&=\frac{\sin t}{\ch \tau}((s\sh R)_s'-\th \tau (s\ch R)_s'),


&\quad Y_t&=\frac{s\sh R}{\ch \tau}(\cos t-\th \tau T_t'),\\


Z_s&=-\th \tau ((s\sh R)_s'-\th \tau (s\ch R)_s'),


&\quad Z_t&=-\frac{s\sh RT_t'}{\ch^2\tau}.


\end{alignat*}


Отсюда находятся выражения для коэффициентов первой квадратичной формы.


Проверим, что средний коэффициент $b_{12}$ второй квадратичной формы равен


нулю. Для этого воспользуемся формулой (43.4) из ([5], с.\,180):


$$


Z_{,ts}=-\Gamma^3_{11}X_sX_t-\Gamma^3_{22}Y_sY_t-


\Gamma^3_{33}Z_sZ_t+b_{12}n^3=


-\frac{1}{Z}X_sX_t-\frac{1}{Z}Y_sY_t+\frac{1}{Z}Z_sZ_t+b_{12}n^3,


$$


где $\Gamma^i_{jk}$ --- символы Кристоффеля метрики $H^3$.


Поскольку $g_{12}=0$, то условие $b_{12}=0$ равносильно следующему:


$Z_{,ts}=\frac{2Z_sZ_t}{Z}$. Смешанная ковариантная производная функции


$Z(s,t)$ на поверхности


$$


Z_{,ts}=\frac{\partial^2Z}{\partial t\partial s}-


\ov \Gamma^1_{12}Z_s-\ov \Gamma^2_{12}Z_t,


$$


где


$\ov \Gamma^1_{12}=\frac{1}{2g_{11}}\frac{\partial g_{11}}{\partial t}$,


$\ov \Gamma^2_{12}=\frac{1}{2g_{22}}\frac{\partial g_{22}}{\partial s}$ ---


символы Кристоффеля внутренней метрики поверхности ([5], c.\,60).


Вычисляем их и смешанную производную от $Z$:


\begin{align*}


\ov \Gamma^1_{12}&=-\frac{sT'_t}{\ch^2\tau ((s\sh R)'_s-


\th \tau (s\ch R)'_s)(s\ch R-\th \tau s\sh R)},\\


\ov \Gamma^2_{12}&=\frac{\ch R((s\sh R)'_s-(s\ch R)'_s \th \tau)}


{\sh R(s\ch R-s\sh R\th \tau)},\\


\frac{\partial^2Z}{\partial t\partial s}&=


- \frac{T'_t((s\sh R)'-2(s\ch R)'_s\th \tau )}{\ch^2\tau}.


\end{align*}


Далее прямым вычислением проверяем, что $Z_{,ts}=\frac{2Z_sZ_t}{Z}$.


\end{pf}





\begin{thm}


Поверхность в $H^3$, заданная уравнениями $(3)$ с произвольными


функциями $R=R(s)>0$, $T=T(t)$, есть поверхность Иоахимсталя.


\end{thm}





{\bf Доказательство.} Согласно лемме 2 имеем


$g_{12}=b_{12}=0$, следовательно,


координатные линии являются линиями кривизны. Линия $t=t_0$ лежит во


вполне геодезической плоскости $X\sin t_0-Y\cos t_0=0$, которая содержит


геодезическую $x=y=0$. Линия $s=s_0$ лежит в некоторой малой 2-сфере


в $H^3$, т.\,к.


$$


X^2(s_0,t)+Y^2(s_0,t)+(Z(s_0,t)-s_0\ch R(s_0))^2=(s_0\sh R(s_0))^2.


\quad\square


$$





{\bf 2.} В [2] доказано, что всякая минимальная поверхность в


$H^3$ c одним семейством плоских линий кривизны является поверхностью


вращения. Покажем, что в классе поверхностей постоянной средней


кривизны без омбилических точек одним семейством плоских линий


кривизны обладают только поверхности вращения (гиперболические


поверхности Делоне). Для этой цели удобно использовать модель геометрии


Лобачевского на гиперболоиде в лоренцевом пространстве $R^{3,1}$, где


$$


H^3=\{x=(x_0,x_1,x_2,x_3) \in R^{3,1}: |x|^2=


-x^2_0+x^2_1+x^2_2+x^2_3=-1,\ x_0>0\}.


$$





Пусть $X(u,v)=\{X^i(u,v),\ i=0,1,2,3\}$ --- поверхность постоянной средней


кривизны $H$ без омбилических точек, вложенная в $H^3$ конформно и


параметризованная своими линиями кривизны. Тогда, как показано в ([2],


теорема 2.1), ее фундаментальными формами являются


$$


ds^2=e^{2 \omega (u,v)}(du^2+dv^2),\quad


\mathrm{II}=(e^{2 \omega}H+1)du^2+(e^{2 \omega}H-1)dv^2.


$$





Уравнение Гаусса имеет вид ([2], теорема 2.1, c.\,6)


\begin{equation}


\Delta \omega +(H^2-1)e^{2 \omega} -e^{-2 \omega}=0.


\end{equation}


Дифференциал Хопфа $\phi (w)dw^2=( \frac{L-N}{2}-iM)dw^2$, где $w=u+iv$,


в данном случае равен $dw^2$ и является голоморфным, что, как


известно, равносильно выполнению условий Петерсона--Кодацци. Поэтому


всякому решению $\omega (u,v)$ уравнения Гаусса в односвязной области


плоскости комплексного переменного $w$ соответствует единственная


поверхность постоянной средней кривизны без омбилических точек.


Деривационные формулы имеют вид


\begin{align*}


X_{uu}&=\omega_uX_u-\omega_vX_v+(e^{2\omega}H+1)n+e^{2\omega}X,\\


X_{uv}&=\omega_vX_u+\omega_uX_v,\\


X_{vv}&=-\omega_uX_u+\omega_vX_v+(e^{2\omega}H-1)n+e^{2\omega}X,\\


n_u&=-(H+e^{-2\omega})X_u,\quad n_v=(-H+e^{-2\omega})X_v,


\end{align*}


где $n(u,v)$ --- вектор единичной нормали к поверхности. В данном репере


векторы $X_u$, $X_v$, $n$ пространственноподобны,


а вектор $X$ времениподобен.





Времениподобность вектора $X=(x_0,x_1,x_2,x_3)$ следует прямо


из определения $H^3$, вложенного в $R^{3,1}$ с метрикой


$\la x,y\ra=-x_0y_0+x_1y_1+x_2y_2+x_3y_3$.


Касательное пространство в точке $x\in H^3$ состоит из векторов


$y\in R^{3,1}$, для которых $\la x,y\ra=0$. Покажем, что для всякого


ненулевого вектора $y\in T_x H^3$ выполнено условие $\la y,y\ra>0$,


откуда и будет следовать пространственноподобность


векторов $X_u$, $X_v$, $n$. Поскольку $x_0>0$ и $\la y,x\ra=0$, то


$y_0=\frac{1}{x_0}\sum\limits^3_{i=1}x_i y_i$. Так как


$x\in H^3$, то $x^2_0=1+\sum\limits^3_{i=1}x^2_i$.


Следовательно,


\begin{equation}


\la y,y\ra=-y^2_0+\sum^3_{i=1}y^2_i=\frac{1}{x^2_0}\bigg(


-\bigg(\,\sum^3_{i=1}x_i y_i\bigg)^2+\bigg(1+


\sum^3_{i=1}x^2_i\bigg)\sum^3_{i=1}y^2_i\bigg)\ge


\frac{\sum\limits^3_{i=1}y^2_i}{x^2_0}\ge0.\tag{$*$}


\end{equation}





В действительности $\la y,y\ra>0$ для всякого ненулевого $y\in T_x H^3$ и


всякого $x\in H^3$. Для того чтобы в первом неравенстве в $(*)$


имело место равенство, необходимо, чтобы $x_i=\lambda y_i$ ($i=1,2,3$),


а второе неравенство в $(*)$ обращается в равенство


только в точке $x=(1,0,0,0)$. Но тогда из условия


$\la y,x\ra=0$ следует, что $y_0=0$, т.\,е. $y=0$.





\begin{thm}


Поверхностями постоянной средней кривизны в $H^3$ без омбилических


точек, имеющими одно семейство линий кривизны, расположенных на вполне


геодезических плоскостях, являются поверхности вращения $($гиперболические


поверхности Делоне$)$.


\end{thm}





\begin{pf}


Условие принадлежности линии кривизны поверхности $X(u,v_0)$ некоторой вполне


геодезической плоскости имеет вид


\begin{equation}


\la X(u,v_0),a(v_0)\ra=-X^0a_0+X^1a_1+X^2a_2+X^3a_3=0,


\end{equation}


где $a(v)$ --- некоторый единичный пространственноподобный вектор:


$\la a,a\ra=1$. Дифференцируя (5) по $u$, получим


$\la X_u,a(v)\ra=0$. Следовательно, разложение


вектора $a(v)$ по подвижному реперу $(X,X_u,X_v,n)$ в $R^{3,1}$ должно иметь


вид $a=c_1X_v+c_2n$. Дифференцирование (5) еще раз по $u$


приводит к соотношению $\la X_{uu},a\ra=0$.


Применяя деривационные формулы, получим


$$


\la\omega_uX_u-\omega_vX_v+(e^{2\omega}H+1)n+e^{2\omega}X,c_1X_v+c_2n\ra=


-\omega_v|X_v|^2c_1+(e^{2\omega}+1)c_2=0.


$$


Отсюда следует


$\frac{c_1}{c_2}=\frac{e^{2\omega}H+1}{e^{2\omega}\omega_v}$.


Так как вектор $a$ единичный, то $|a|^2=c_1^2e^{2\omega}+c_2^2=1$.


Находим разложение вектора $a$ по подвижному реперу


$X$, $X_u$, $X_v$, $n$:


$$


a(v)=\frac{(e^{2\omega}H+1)X_v+e^{2\omega}\omega_vn}


{e^{\omega}((e^{2\omega}H+1)^2 +e^{2\omega}\omega_v^2)^{1/2}}.


$$


Продифференцируем это выражение по $u$ и еще раз используем деривационные


формулы. Получим


\begin{multline*}


\bigg(L_uX_v+L(\omega_vX_u+\omega_uX_v)+(e^{2\omega}\omega_v)_un+


e^{2\omega}\omega_v\bigg(-\frac{L}{e^{2\omega}}X_u\bigg)\bigg)


(L^2+e^{2\omega}\omega_v^2)-\\


-(LX_v+e^{2\omega}\omega_vn)(\omega_u(L^2+e^{2\omega}\omega_v^2)+LL_u+


e^{2\omega}\omega_v(\omega_u\omega_v+\omega_{uv}))=0,


\end{multline*}


где $L=He^{2\omega}+1$. Приравнивая нулю коэффициенты при базисных векторах


$X_v$ и $n$, получим систему из двух дифференциальных уравнений,


которым должна удовлетворять функция $\omega (u,v)$,


\begin{align*}


(L_u+L\omega_u)(L^2+e^{2\omega}\omega_v^2)-L(\omega_uL^2+


2e^{2\omega}\omega_u\omega_v^2+LL_u+e^{2\omega}\omega_v\omega_{uv})&=0,\\


(e^{2\omega}\omega_v)_u(L^2+e^{2\omega}\omega^2_v)-e^{2\omega}\omega_v


(\omega_uL^2+2e^{2\omega}\omega_u\omega_v^2+LL_u+e^{2\omega}\omega_v


\omega_{uv})&=0.


\end{align*}


Эта система приводится к одному уравнению


$$


\omega_{uv}(He^{2\omega}+1)-\omega_u\omega_v(He^{2\omega}-1)=0.


$$


Если $\omega_u \ne0$, то уравнение интегрируется, т.\,к.


$$


(\log \omega_u)_v=\frac{\omega_{uv}}{\omega_u}=\frac{He^{\omega}-e^{-\omega}}


{He^{\omega}+e^{-\omega}}\omega_v=(\log (He^{\omega}+e^ {-\omega}))_v.


$$


Следовательно,


\begin{equation}


\omega_u=C(u)(He^{\omega}+e^{-\omega}).


\end{equation}





Теперь воспользуемся теоремой, доказанной Г.Венте.


\medskip





{\bf Теорема} ([2], теорема 2.3).


{\it Пусть $\omega (u,v)$ есть решение системы


дифференциальных уравнений


\begin{equation}


\begin{split}


&\Delta \omega +Ae^{2\omega}-Be^{-2\omega}=0,\\


&2\omega_u=\alpha (u)e^{\omega}+\beta (u)e^{-\omega}.


\end{split}


\end{equation}


Если $\omega_v \ne0$, то функции $\alpha (u)$, $\beta (u)$ суть


решения системы уравнений


\begin{equation}


\begin{split}


\alpha ''&=a \alpha -2\alpha^2 \beta -2A\beta ,\\


\beta ''&=a\beta -2\alpha \beta^2-2B\alpha,


\end{split}


\end{equation}


где $a$ есть некоторая постоянная.}


\medskip





В условиях доказываемой теоремы первое из уравнений системы (7)


совпадает с уравнением Гаусса (4), причем $A=H^2-1$, $B=1$.


Второе из уравнений системы (7) есть уравнение (6), причем


$\alpha (u)=2HC(u)$, $\beta (u)=2C(u)$. По теореме


Венте функция $C(u)$ должна удовлетворять системе уравнений (8), которая имеет


вид


\begin{equation}


\begin{split}


2HC''&=2aHC-16H^2C^3-4(H^2-1)C,\\


2C''&=2aC-16HC^3-4HC.


\end{split}


\end{equation}





1) Если $H \ne0$, то, умножая на $H$ второе уравнение системы (9) и сравнивая


его с первым уравнением, находим $4H^2C=4(H^2-1)C$.


Следовательно, $C\equiv 0$.





2) Если $H=0$, то $\alpha =0$, и из первого уравнения следует, что $C\equiv0$.





Следовательно, $\omega_u\equiv 0$ и функция $\omega$ зависит только от


переменной $v$. Как следует из доказанной ниже леммы, поверхность $X(u,v)$


будет поверхностью вращения в пространстве Лобачевского.


\end{pf}





Как уже отмечалось, для минимальных поверхностей в $H^3$ c плоским


семейством линий кривизны теорема была доказана Г.\,Венте


([2], теорема 5.5).


Следующая лемма справедлива для поверхностей, у которых первая и вторая


фундаментальные формы зависят от одной переменной, а средняя кривизна $H$ не


обязательно постоянна.





\begin{lem}


Пусть первая и вторая фундаментальные формы поверхности $X(u,v)$


в пространстве Лобачевского зависят от одной переменной $u$ и имеют вид


\begin{align*}


\mathrm I&=e^{2 \omega (u)}(du^2+dv^2),\\


\mathrm{II}&=(e^{2 \omega (u)}H+1)du^2+(e^{2 \omega (u)}H-1)dv^2.


\end{align*}


Тогда поверхность $X(u,v)$ является поверхностью вращения в $H^3$.


\end{lem}





\begin{pf}


Покажем, что линия $X(u_0,v)$ имеет постоянную кривизну


и нулевое кручение в $H^3$, т.\,е. лежит в некоторой вполне


геодезической плоскости. Единичный касательный вектор к линии $X(u_0,v)$


есть $\tau =e^{-\omega (u_0)}X_v(u_0,v)$. Вычислим его ковариантную


производную в $H^3$, которая по формуле Френе будет равна $k_1\nu$,


где $k_1$ --- кривизна линии, а $\nu$ --- вектор главной нормали.


Используя выписанные выше деривационные формулы, с учетом того, что


$\omega_v=0$, получим


\begin{multline*}


\frac{D\tau}{ds}=\PR_{T_XH^3}\frac{d\tau}{ds}=\PR_{T_XH^3}e^{-\omega(u_0)}


\frac{d\tau}{dv}=\PR_{T_XH^3}e^{-2\omega(u_0)}X_{vv}=\\


=e^{-2\omega}(-\omega 'X_u+(e^{2\omega}H-1)n)=k_1\nu.


\end{multline*}


Отсюда видно, что $k_1^2=|\frac{D\tau}{ds}|^2=e^{-4\omega


(u_0)}(\omega'{}^2 e^{2\omega}+(e^{2\omega}H-1)^2)$ не зависит от


переменной $v$ вдоль линии $X(u_0,v)$. Cледовательно, кривизна этой


линии постоянна. Вектор $\beta $ бинормали кривой $X(u_0,v)$


перпендикулярен в $R^{3,1}$ к $\tau$, $\nu$, $X$. Поскольку


$\nu =\lambda(u_0)(-\omega 'X_u+(e^{2\omega}H-1)n)$


и


$|X_u|=e^{\omega(u_0)}$, то


ясно, что $\beta =\lambda (u_0)e^{-\omega (u_0)}((e^{2\omega}H-1)X_u+


\omega'e^{2\omega}n)$. Вычислим ковариантную производную


$\dfrac{D\beta}{ds}$


вектора бинормали вдоль кривой $X(u_0,v)$, которая по формулам


Френе равна $-k_2\nu $,


\begin{multline*}


\frac{D \beta}{ds}=\PR_{T_XH^3}\frac{d \beta}{ds}=\PR_{T_XH^3}e^{- \omega}


\frac{d \beta}{dv}=


\PR_{T_XH^3} \lambda e^{-2 \omega}((e^{2 \omega}H-1)X_{uv}+\omega '


e^{2 \omega}n'_v)=\\


=\PR_{T_XH^3}\bigg((e^{2\omega}H-1)\omega 'X_v+\omega 'e^{2\omega}


\bigg(-\frac{e^{2\omega}H-1}{e^{2\omega}}\bigg)X_v\bigg)=0.


\end{multline*}


Следовательно, кручение $k_2$ линии $X(u_0,v)$ равно нулю,


и эта линия плоская.





Одна из деривационных формул $X_{uv}=\omega 'X_v$ легко интегрируется:


$$


X(u,v)=e^{\omega (u)} \Phi (v)+ \Psi (u),


$$


где $\Phi (v)=(\phi_0,\phi_1,\phi_2,\phi_3)$,


$\Psi(u)=(\psi_0,\psi_1,\psi_2,\psi_3)$ --- некоторые


вектор-функции из $R^{3,1}$.


Так как $|X'_v|=e^{\omega}$, то $| \Phi '(v)|=1$.


Поскольку линия $X(u_0,v)$ плоская, то можно выбрать базис в $R^{3,1}$


таким образом, чтобы $X^3(u_0,v) \equiv 0$. Тогда


$X(u_0,v) \subset H^2_0=\{x_0,x_1,x_2,0) \mid x_0^2+x_1^2+x_2^2=-1\}$.


Так как кривизна линии $X(u_0,v)$ постоянна, то это либо


1)~окружность, либо 2)~геодезическая, либо


3)~орицикл, либо 4)~``псевдоокружность'' (т.\,е. линия, которая в модели


Пуанкаре для $H^2$ в верхней полуплоскости представляется дугой


окружности, расположенной над осью $Ох$). Действительно, в этой модели


для кривой, параметризованной длиной дуги, имеем $x'{}^2+y'{}^2=y^2$, орт


касательной $\tau =(x',y')$, орт главной нормали равен $\nu =(-y',x')$.


В случае постоянной кривизны $k$ кривой формула Френе


$\frac{D\tau}{dt}=k\nu$ принимает вид


\begin{align*}


x''-\frac{2}{y}x'y'&=-ky',\\


y''+\frac{x'{}^2-y'{}^2}{y}&=kx'.


\end{align*}


Система интегрируется так же, как в ([6], с.\,80). В результате получаем


$(x{-}c_1)^2{+}(\frac{k}{c}{+}y)^2{=}c^{-2}$. В зависимости от величины $k$ и


возникают перечисленные возможности. Последовательно разберем их на


модели $H^2_0 \subset R^{2,1}$. 1)~Если кривая $X(u_0,v)$ представляет


собой окружность, то после поворота системы координат в $R^{3,1}$ ее


уравнения примут вид


$X(u_0,v)=(a,\,\sqrt{a^2-1}\cos v$, $\sqrt{a^2-1}\sin v,\,0)$.


%%%%%% join in eng.              ^^^^


Тогда $\Phi '=(0,-\sin v,\cos v,0)$. За счет сдвига начала координат


можно считать, что $\Phi=(0,\cos v,\sin v,0)$, и условие $\la X_u,X_v\ra=0$


приводит к тому, что $\la\Phi ',\Psi '\ra=0$ и, значит, для всех $v$ должно


выполняться уравнение $-\sin v\psi_1'+\cos v\psi_2'=0$, откуда


$\Psi=(\psi_0,0,0,\psi_3)$. Следовательно, уравнения поверхности имеют вид


$X(u,v)=(\psi_0(u),\,e^{\omega (u)}\cos v,\,e^{\omega (u)}\sin v,\,


\psi_3(u))$, т.\,е. поверхность инвариантна относительно группы вращений,


изоморфной $SO(2)$. 2)~Если $X(u_0,v)$ --- ``псевдоокружность'' или


геодезическая, то за счет поворота системы координат в $R^{3,1}$ можно


считать, что она имеет уравнения


$X(u_0,v)=(a\ch v,\, a\sh v,\,\sqrt{a^2-1},\,0)$,


где $a>1$ в случае ``псевдоокружности''


и $a=1$ в случае геодезической. Так как


$|\Phi '|=1$, то $\Phi '=(\sh v,\ch v,0,0)$, и за счет выбора начала


координат можно считать, что $\Phi =(\ch v,\sh v,0,0)$. Затем условие


$\la X_u,X_v\ra=0$ дает $\Psi=(0,0,\psi_2,\psi_3)$,


и уравнения поверхности принимают вид


$X(u,v)=(e^{\omega (u)}\sh v,\,e^{\omega (u)}\ch v,\,


\psi_2(u),\,\psi_3(u))$.


Эта поверхность инвариантна относительно подгруппы


собственных вращений плоскости Лоренца.


3)~Остается последняя возможность, когда при любом $u_0$ кривая $X(u_0,v)$


есть орицикл. Зафиксируем значение $u_0$ и рассмотрим данную кривую. За счет


вращения системы координат в $R^{3,1}$ можно добиться, чтобы ее уравнение


приняло вид


$X(u_0,v)=e^{\omega (u_0)}\Phi (v)+\Psi (u_0)=


\big(\dfrac{v^2}{2}+1,v,\dfrac{v^2}{2},0\big)$.


Следовательно, $X'_v(u_0,v)=e^{\omega (u_0)}\Phi '(v)=(v,1,v,0)$. Tак как


$|\Phi '|=1$, то $e^{\omega (u_0)}=1$, и $\Phi '=(v,1,v,0)$. Интегрируя это


уравнение и подходящим образом выбирая начало координат, получим


$\Phi (v)=\big(\frac{v^2}{2}+1,v,\frac{v^2}{2},0\big)$. Далее, условие


$\la X_u,X_v\ra=0$ сводится к $-\psi_0'v+\psi_1'+\psi_2'v=0$. Отсюда находим


$\psi_1=c_1$, $\psi_2-\psi_0=c_2$. Значит,


$\Psi =(\psi_0(u),c_1,\psi_0+c_2,\psi_3(u))$,


и поверхность имеет вид


$$


X(u,v)=e^{\omega (u)}\bigg(\dfrac{v^2}{2}+1,v,\dfrac{v^2}{2},0\bigg)+


(\psi_0(u),c_1,\psi_0(u)+c_2,\psi_3(u)).


$$


Поскольку при всяком фиксированном $u_0$ линия $X(u_0,v)$ есть орицикл, то


$e^{\omega (u)} \equiv 1$ (т.\,к. кривизна всякого орицикла равна единице).


Теперь надо воспользоваться тем, что поверхность


$X(u,v)=\big(\frac{v^2}{2}+1+\psi_0,v+c_1,\frac{v^2}{2}+


\psi_0+c_2,\psi_3\big)$


удовлетворяет условию $|X(u,v)|^2=-1$. Записав это условие, получим


$2\psi_0(c_2-1)+c_2v^2+2vc_1+c_1^2+c_2^2+\psi_3^2\equiv 0$. Отсюда следует,


что 1)~$c_1=c_2=0$, 2)~$\psi_3^2=2\psi_0$. Cледовательно, если ввести


новую переменную $u_1=\psi_3(u)$, то уравнения поверхности примут вид


$X(u,v)=\big(\frac{v^2}{2}+1+\frac{u_1^2}{2},v,\frac{v^2}{2}+


\frac{u_1^2}{2},u_1\big)$.


Поэтому в третьем случае поверхность представляет собой область на орисфере.


\end{pf}





\begin{thebibliography}{9}





\bibitem{1}


Bianchi L. {\em Lezioni di geometria differenziale}. -- Bologna: Edit. Nic.


Zanicheli, 1927. -- V.\,1. -- 801~p.


\bibitem{2}


Wente H.C. {\em Constant mean curvature immersions of Enneper type} //


Memoirs Amer. Math. Soc. -- 1992. -- \No\,478. -- P.\,1--77.


\bibitem{3}


Шуликовский В.И. {\em Классическая дифференциальная геометрия в тензорном


изложении}. -- М.: Физматлит, 1963. -- 540~с.


\bibitem{4}


Бердон А. {\em Геометрия дискретных групп}. -- М.: Наука,


1986. -- 300~с.


\bibitem{5}


Эйзенхарт Л.П. {\em Риманова геометрия}. -- M.: Ин. лит., 1948. -- 316 с.


\bibitem{6}


Трофимов В.В. {\em Введение в геометрию многообразий с симметриями}. -- М.:


Изд-во МГУ, 1989. -- 353~с.








\end{thebibliography}





\vskip 2cm





\post{\it Харьковский государственный университет}{\it Поступила}


\post{}{08.06.1999}





\label{end}


\end{document}


